\documentclass{standalone}
\usepackage{circuitikz}
\usetikzlibrary{calc}
\definecolor{circuitnotemark}{HTML}{FE00FE}
\begin{document}
\begin{circuitikz}[american, line width=0.8pt]
\ctikzset{bipoles/length=1.2cm}
\ctikzset{grounds/scale=1.36}
\foreach \x in {0,2,4,6,8,10,12,14,16} {\foreach \y in {0,-2,-4} {\fill[gray, opacity=0.35] (\x,\y) circle (1.1pt);}}
\node[gray, font=\scriptsize] at (0,0.84) {1};
\node[gray, font=\scriptsize] at (2,0.84) {2};
\node[gray, font=\scriptsize] at (4,0.84) {3};
\node[gray, font=\scriptsize] at (6,0.84) {4};
\node[gray, font=\scriptsize] at (8,0.84) {5};
\node[gray, font=\scriptsize] at (10,0.84) {6};
\node[gray, font=\scriptsize] at (12,0.84) {7};
\node[gray, font=\scriptsize] at (14,0.84) {8};
\node[gray, font=\scriptsize] at (16,0.84) {9};
\node[gray, font=\scriptsize, anchor=east] at (-0.84,0) {a};
\node[gray, font=\scriptsize, anchor=east] at (-0.84,-2) {b};
\node[gray, font=\scriptsize, anchor=east] at (-0.84,-4) {c};
\coordinate (a5) at (8,0);
\coordinate (b5) at (8,-2);
\coordinate (c5) at (8,-4);
\coordinate (c7) at (12,-4);
\coordinate (c9) at (16,-4);
\coordinate (c3) at (4,-4);
\draw (a5) node[ocirc]{} node[above left]{IN}; % line 3
\draw (a5) to[R, l_=$R_{1}$, a^=$10\,\mathrm{k}\Omega$] (b5); % line 4
\node[npn, rotate=-90] (part-Q1) at (c5) {}; % line 5
\draw (c7) to[R, l_=$R_{C}$, a^=$1\,\mathrm{k}\Omega$] (c9); % line 6
\node[vcc] at (c9) {VCC}; % line 7
\node[ground, rotate=-90] at (c3) {}; % line 8
\draw (b5) -- (part-Q1.base); % line 10
\draw (part-Q1.collector) -- (c7); % line 11
\draw (part-Q1.emitter) -- (c3); % line 12
\path (20,-8.869) rectangle (27.959,0.593); % line 14
\node[anchor=west, inner sep=0, circuitnotemark, font=\large] at (20,0) {X}; % line 14
\node[anchor=west, inner sep=0, circuitnotemark, font=\large] at (20,-0.487) {X}; % line 14
\node[anchor=west, inner sep=0, circuitnotemark, font=\large] at (20,-0.974) {X}; % line 14
\node[anchor=west, inner sep=0, circuitnotemark, font=\large] at (20,-1.46) {X}; % line 14
\node[anchor=west, inner sep=0, circuitnotemark, font=\large] at (20,-1.947) {X}; % line 14
\node[anchor=west, inner sep=0, circuitnotemark, font=\large] at (20,-2.434) {X}; % line 14
\node[anchor=west, inner sep=0, circuitnotemark, font=\large] at (20,-2.921) {X}; % line 14
\node[anchor=west, inner sep=0, circuitnotemark, font=\large] at (20,-3.408) {X}; % line 14
\node[anchor=west, inner sep=0, circuitnotemark, font=\large] at (20,-3.895) {X}; % line 14
\node[anchor=west, inner sep=0, circuitnotemark, font=\large] at (20,-4.381) {X}; % line 14
\node[anchor=west, inner sep=0, circuitnotemark, font=\large] at (20,-4.868) {X}; % line 14
\node[anchor=west, inner sep=0, circuitnotemark, font=\large] at (20,-5.355) {X}; % line 14
\node[anchor=west, inner sep=0, circuitnotemark, font=\large] at (20,-5.842) {X}; % line 14
\node[anchor=west, inner sep=0, circuitnotemark, font=\large] at (20,-6.329) {X}; % line 14
\node[anchor=west, inner sep=0, circuitnotemark, font=\large] at (20,-6.816) {X}; % line 14
\node[anchor=west, inner sep=0, circuitnotemark, font=\large] at (20,-7.302) {X}; % line 14
\node[anchor=west, inner sep=0, circuitnotemark, font=\large] at (20,-7.789) {X}; % line 14
\node[anchor=west, inner sep=0, circuitnotemark, font=\large] at (20,-8.276) {X}; % line 14
\node[anchor=south west, inner sep=2pt, circuitnotemark, font=\LARGE\bfseries] (circuittitle) at (current bounding box.north west) {X};
\path (circuittitle.south west) rectangle ++(9.418,0.853);
\end{circuitikz}
\end{document}
